% Generated draft: requires mathematical and source review.
\documentclass[11pt]{article}
\usepackage[margin=25mm]{geometry}
\usepackage{fontspec}
\setmainfont{Latin Modern Roman}
\usepackage{amsmath,amssymb,mathtools,unicode-math}
\setmathfont{Latin Modern Math}
\usepackage{enumitem,needspace,xurl,hyperref}
\setlength{\parindent}{0pt}
\setlength{\parskip}{5pt}
\setlength{\emergencystretch}{3em}
\newcommand{\N}{\mathbb{N}}
\newcommand{\Z}{\mathbb{Z}}
\newcommand{\Q}{\mathbb{Q}}
\newcommand{\R}{\mathbb{R}}
\newcommand{\C}{\mathbb{C}}
\newcommand{\F}{\mathbb{F}}
\newcommand{\A}{\mathbb{A}}
\newcommand{\PP}{\mathbb{P}}
\newcommand{\E}{\mathbb{E}}
\newcommand{\eps}{\varepsilon}
\newcommand{\dd}{\,\mathrm{d}}
\begin{document}
\section*{UnsolvedMath ID: 800408}
\textbf{Source problem code:} catalogue entry
\subsection{Definitions and mathematical statement}
Let \(N,D,d\) be integers with \(N\ge 1\), \(D\ge 1\), and \(d\ge 2\). For a number field \(K\) and a degree-\(d\) morphism \(f:\PP^N\to\PP^N\) defined over \(K\), write \(f^{\circ 0}=\operatorname{id}\) and define the forward orbit of \(P\in\PP^N(K)\) by
\[
\mathcal O_f(P)=\{f^{\circ n}(P):n\in\Z_{\ge 0}\}.
\]
The orbit is Zariski dense if its Zariski closure in \(\PP^N_{\overline K}\) is all of \(\PP^N_{\overline K}\). Let \(h\) be the absolute logarithmic Weil height on projective space associated with \(\mathcal O_{\PP^N}(1)\). The canonical height is
\[
\widehat h_f(P)=\lim_{n\to\infty}\frac{h(f^{\circ n}(P))}{d^n}.
\]
The notation \(h_{\mathcal M}(f)\) denotes the moduli height appearing in the supplied source statement, intended as a height of the conjugacy class of \(f\) under projective linear changes of coordinates. The supplied record does not specify the moduli space, height construction, or normalization. These choices must be recovered from the source before this term is fully defined; no particular choice is imposed here.

\textbf{Dynamical Lang Height Conjecture (supplied formulation).} There exist positive real constants \(c_1(D,N,d)\) and \(c_2(D,N,d)\), depending only on \(D,N,d\), such that for every number field \(K/\Q\) with \([K:\Q]\le D\), every degree-\(d\) morphism \(f:\PP^N\to\PP^N\) defined over \(K\), and every point \(P\in\PP^N(K)\) with Zariski-dense forward orbit, one has
\[
\widehat h_f(P)\ge c_1(D,N,d)\,h_{\mathcal M}(f)+c_2(D,N,d).
\]
The positive additive term is retained exactly as supplied. Its formulation requires the intended choice and normalization of \(h_{\mathcal M}\) to be clarified.
\subsection{Short English statement}
The supplied conjecture asserts a uniform positive lower bound for the canonical height of a rational point with Zariski-dense orbit, growing linearly with the moduli height of the map. The constants depend only on the field-degree bound, dimension, and map degree. The moduli height is not defined in the supplied record, so its precise interpretation requires review.
\subsection{Sources}
\begin{itemize}
\item Local UnsolvedMath registry, permanent ID \(800408\), source problem code \texttt{catalogue entry}, titled ``Open Problems in Arithmetic Dynamics --- Conjecture 31 (Dynamical Lang Height Conjecture).'' The registry labels this record \emph{open}; this is supplied metadata, not an independently verified status assessment.
\item Benedetto, DeMarco, Ingram, Jones et al., \emph{Current Trends and Open Problems in Arithmetic Dynamics} (2018), Conjecture 31. Supplied source link: \url{https://arxiv.org/abs/1806.04980}. The record attributes the problem to these authors and gives the context ``Dynamical Lang Height Conjecture.''
\item Supplied provenance: extraction method \texttt{arxiv\_latex}; source access date \texttt{2026-10-16}. This access information is supplied provenance and does not indicate an independent source visit or verification for this draft.
\end{itemize}
\end{document}
